This paper studies a generalization failure discovered during an automated architecture search for NanoGPT, which we call \emph{forking}. When the training data is repeatedly consumed, models with an over-encoding n-gram memory branch exhibit a clear separation between training and validation loss, and this separation often grows in a staircase pattern aligned with replay epoch boundaries. We study this phenomenon in a controlled vanilla NanoGPT setting. Our experiments show that over-encoding partitions the optimization process into finer context-specific subspaces. Repeated updates then sharpen the observed continuations within these subspaces while damaging alternative continuations that remain plausible but receive insufficient coverage. We further observe the same forking phenomenon in a DeepSeek-style model with an Engram-like memory module. In a continued-pretraining stress test on small ID and OOD corpora, the gap reopens when the n-gram table remains trainable; freezing the table substantially reduces but does not eliminate the effect.
